\documentclass[exercice,a4]{thedocument}%
\usepackage[french]{babel}%
\usepackage{pstricks}%
\usepackage{pst-plot}%
\usepackage{multido}%
\startdatakeys
    \source{22GENMATMEAG2}
    \BO{2026}
    \cycle{4}
    \competencesDuSocle{CA,RA}
    \objectifsApprentissage{C4C2aa,C4D3ac,C4D3ba}
    \motsCles{lecture graphique,fonction linéaire,vitesse,durée}
\enddatakeys
\begin{document}%

\medskip

Le Tour de France 2021 est une compétition de cyclisme s'étant déroulée du 26 juin au 18 juillet 2021.

L'étape 5 était un contre-la-montre individuel de $27,2$~km. Le graphique ci-dessous indique le temps effectué par Tadej Pogacar en fonction de la distance parcourue sur son trajet (source : https://www.letour.fr).

\medskip

On répondra aux questions suivantes à l'aide de ce graphique reproduit en annexe.

On laissera les traits de construction éventuels sur le graphique de l'annexe.

\begin{center}
\psset{xunit=0.2cm,yunit=0.15cm,arrowsize=2pt 3}
\begin{pspicture}(-5,-5)(32,43)
\psaxes[linewidth=1.25pt,labelFontSize=\scriptstyle,Dx=5,Dy=5]{->}(0,0)(0,0)(32,43)
\multido{\n=0+1}{33}{\psline[linewidth=0.1pt](\n,0)(\n,43)}
\multido{\n=0+5}{7}{\psline[linewidth=0.5pt](\n,0)(\n,43)}
\multido{\n=0+1}{44}{\psline[linewidth=0.1pt](0,\n)(32,\n)}
\multido{\n=0+5}{9}{\psline[linewidth=0.5pt](0,\n)(32,\n)}
\psline[linewidth=1.25pt,linecolor=blue](27.2,32)
\uput[u](26,0){distance (en km)}\uput[r](0,42){durée (en minutes)}
\end{pspicture}
\end{center}

\begin{enumerate}
\item Quelle distance approximative Tadej POGACAR a-t-il parcouru en 20 minutes ?
\item En combien de temps environ Tadej POGACAR a-t-il effectué les $27,2$ km du parcours ?
\item Le cycliste Bryan Coquard a lui aussi effectué ce contre-la-montre.

Le temps, en minutes, mis par Bryan Coquard pour parcourir la distance $x$, exprimée en km, peut être représenté par la fonction linéaire $f$ d'expression algébrique:

\[f(x) = 1,4x.\]

	\begin{enumerate}
		\item Représenter la fonction $f$ sur le graphique en annexe.
		\item Lequel de ces deux coureurs a été classé avant l'autre ?
	\end{enumerate}
\item Est-il vrai que la vitesse moyenne de Tadej Pogacar sur l'ensemble du parcours est
supérieure à $50$ km/h ? Justifier.
\end{enumerate}

\bigskip


\textbf{Annexe de l'exercice 4}

\medskip

\begin{center}
\psset{xunit=0.3cm,yunit=0.2cm,arrowsize=2pt 3}
\begin{pspicture}(-5,-5)(32,43)
\psaxes[linewidth=1.25pt,labelFontSize=\scriptstyle,Dx=5,Dy=5]{->}(0,0)(0,0)(32,43)
\multido{\n=0+1}{33}{\psline[linewidth=0.1pt](\n,0)(\n,43)}
\multido{\n=0+5}{7}{\psline[linewidth=0.5pt](\n,0)(\n,43)}
\multido{\n=0+1}{44}{\psline[linewidth=0.1pt](0,\n)(32,\n)}
\multido{\n=0+5}{9}{\psline[linewidth=0.5pt](0,\n)(32,\n)}
\psline[linewidth=1.25pt,linecolor=blue](27.2,32)
\uput[u](27,0){distance (en km)}\uput[r](0,42){durée (en minutes)}
\end{pspicture}
\end{center}

\showthesource
\end{document}
